\documentclass{article}

\usepackage{amsmath}
\usepackage{amssymb}
\usepackage{graphicx}
\usepackage{booktabs}
\usepackage{url}
\usepackage[a4paper, margin=1in]{geometry}

\title{Diffusion Model Mathematics Notes}
\author{Rich Wu}
\date{\today}

\begin{document}

\maketitle

\section{Diffusion Model Derivation}
\subsection{Master Equations}
De-nosing Diffusion Probabilistic Model (DDPM) is an application of stochastic Markov process, 
derived from the following continuous time VP-SDE (Variance-Preserving Stochastic Differential Equation):
\begin{equation}
  d\textbf{X}(t)=-\frac{\beta(t)}{2} \textbf{X}(t)dt+\sqrt{\beta(t)}d\textbf{W}(t)
\end{equation}

$\textbf{X}(t)$ is the stochastic random variable (n-dim, in general), 
by definition $d\textbf{X}(t)\equiv \textbf{X}(t+dt)-\textbf{X}(t)$.

$\textbf{W}(t)$ is the standard Wiener process, 
by definition $d\textbf{W}(t)\equiv \textbf{W}(t+dt)-\textbf{W}(t)=\sqrt{dt}\epsilon$, 
where $\epsilon$ is the Gaussian random number sampled from normal distribution with mean=0 and variance=1.

$\beta(t)\geq 0$, for $t \geq 0$, is the variance scheduler function.

\paragraph{}
This SDE is known as the time-dependent \textbf{Ornstein-Uhlenbeck} process in physics, 
where the 1st-term in the right-hand-side of SDE is the drift term, and 2nd-term is the diffusion term.
Physically it was to describe the particle velocity distribution under a drift force and a fluctuation (noise) force.
\paragraph{}
The solution of $\textbf{X}(t)$, or equivalently the distribution function $q(\textbf{X}(t))$,
can be found by solving the corresponding 
\textbf{Fokker-Planck equation} (aka \textbf{Kolmogorov forward equation}),
which is the equation to describe the time evolution of conditional probability $q(x_t|x_s), t > s\geq 0$:
\begin{equation}
  \frac{\partial}{\partial t}q(x_t|x_s)=\frac{\beta(t)}{2} \left( \frac{\partial}{\partial x_t} (x_tq(x_t|x_s))
  + \frac{\partial^2 q(x_t|x_s)}{\partial x_t^2} \right) ; t > s \geq 0
\end{equation}
Here we use the short notation for conditional probability density function:
\begin{equation}
  q(\textbf{X}(t)=x; t=t | \textbf{X}(s)=y; t=s) \equiv q(x;t|y;s) \equiv q(x_t|x_s)
\end{equation}
note that $\textbf{X}(t)$ is the stochastic random variable, 
and $x, y$ are the sampled results at time $t, s$, for $t>s$.

\subsection{Analytic Solution}
The solution can be found by changing the variables:
\begin{eqnarray}
  B(t) &=& \int_0^t \beta(\tau)d\tau \\
  \alpha_t \equiv \alpha(t) &=& \exp(-B(t)) \\
  y_t &=& x_t \exp(B(t)/2)
\end{eqnarray}
The analytic solution is:
\begin{equation}
  q(x_t|x_s)=\exp \left( -\frac{(x_t-\sqrt{\alpha_t / \alpha_s}x_s)^2}{2(1-\alpha_t / \alpha_s)} \right)
\end{equation}
This is the Gaussian distribution, here we drop the normalization constant.

By definition, $\alpha(t=0)=1$, so:
\begin{equation}
  q(x_t|x_0)=\exp \left( -\frac{(x_t-\sqrt{\alpha_t}x_0)^2}{2(1-\alpha_t)} \right)
\end{equation}

\subsubsection{forward sampling}
For the forward sampling (noising process) at arbitrary time $t > 0$:
\begin{equation}
  x_t = \sqrt{\alpha_t}x_0 + \sqrt{1-\alpha_t} \epsilon_t
\end{equation}
where $\epsilon_t$ is the Gaussian noise.
For a given initial data $x_0$ and a pre-defined signal decay function $\alpha(t)$, 
one can compute "noised data" at arbitrary time $t$ using forward sampling formula.

\subsubsection{signal decay function $\alpha_t$}
The signal decay function $\alpha_t$ (called signal power in some other literatures)
controls how the initial data (input image) lost its signature with gradually increasing noise.

Two common types of $\alpha_t$ are:

\textbf{linear scheduler}(original DDPM paper)
\begin{eqnarray}
  \beta(t) &=& \frac{1}{T}\left(\beta_0+ (\beta_1-\beta_0)\frac{t}{T} \right)  ; 0\leq t \leq T  \\
  \alpha_t &=& \exp \left( -\int_0^t \beta(\tau)d\tau \right)
\end{eqnarray}

\textbf{cosine scheduler}(proposed in iDDPM paper):
\begin{equation}
  \alpha_t = \cos^2 \left( \frac{\pi}{2}\frac{t}{T} \right)  ; 0\leq t \leq T
\end{equation}

\begin{figure}[h]
  \centering
  \caption{signal decay function $\alpha_t$}
  \fbox{\includegraphics[width=5cm]{example-image}}
  \label{fig1}
\end{figure}

\subsubsection{reverse sampling}
Ideally if we know the reverse conditional probability $q(x_s|x_t) (s<t)$, 
we are able to apply the reverse sampling in the reverse diffusion process, 
and back to original initial data $x_0$.
However, the conditional probability $q(x_s|x_t)$ is unknown unless the marginal probability $q(x_0)$ 
of the initial data distribution is given. 
Unfortunately, $q(x_0)$ is usually unknown and is what we want to model. 

The key step of DDPM is to use $q(x_s|x_t, x_0)$ to approximate the $q(x_s|x_t)$, 
where $q(x_s|x_t, x_0)$, the probability of sampling $x_s$ at time $s < t$ when conditioned on
$x_t$ and $x_0$, can be calculated analytically by (for $t>s>0$) using Bayers' rule:
\begin{equation}
  q(x_s|x_t,x_0)=\frac{q(x_t, x_s, x_0)}{q(x_t,x_0)}=\frac{q(x_t|x_s,x_0)q(x_s|x_0)q(x_0)}{q(x_t|x_0)q(x_0)}
  = \frac{q(x_t|x_s)q(x_s|x_0)}{q(x_t|x_0)}
\end{equation}
\begin{equation}
  q(x_s|x_t,x_0)=\exp \left(\frac{(x_s-\mu_s(t,x_t,x_0))^2}{2\Sigma_{s,t}^2}   \right)  
\end{equation}
\begin{eqnarray}
  \mu_s(t, x_t, x_0)&=&\frac{1-\alpha_s}{1-\alpha_t}\sqrt{\frac{\alpha_t}{\alpha_s}}x_t
  + \frac{1-\alpha_t/\alpha_s}{1-\alpha_t}\sqrt{\alpha_s}x_0  \\
  \Sigma_{s,t}^2 &=&\frac{1-\alpha_s}{1-\alpha_t} \left( 1-\frac{\alpha_t}{\alpha_s} \right)
\end{eqnarray}

This is also a Gaussian distribution with mean $\mu_s$ and variance $\Sigma^2_{s,t}$,

\subsubsection{DDPM reverse sampling formula}
\begin{eqnarray}
  x_0^\theta(x_t) &=& \frac{x_t-\sqrt{1-\alpha_t} \epsilon_t^\theta(x_t)}{\sqrt{\alpha_t}} \\
  x_s &=& \frac{1-\alpha_s}{1-\alpha_t}\sqrt{\frac{\alpha_t}{\alpha_s}}x_t 
  + \frac{1-\alpha_t/\alpha_s}{1-\alpha_t} \sqrt{\alpha_s} x_0^\theta(x_t) + \Sigma_{s,t} \epsilon
\end{eqnarray}

\subsubsection{DDIM reverse sampling formula}
\begin{eqnarray}
  x_0^\theta(x_t) &=& \frac{x_t-\sqrt{1-\alpha_t} \epsilon_t^\theta(x_t)}{\sqrt{\alpha_t}} \\
  x_s &=& \sqrt{\alpha_s}x_0^\theta(x_t) + \sqrt{1-\alpha_s-\eta \Sigma_{s,t}^2} \epsilon_t^\theta(x_t) +
  \eta \Sigma_{s,t} \epsilon \\
  \eta &=& 1   \text{,        mathematically identical to DDPM sampling} \\
  \eta &=& 0   \text{,        deterministic reverse sampling}
\end{eqnarray}



\end{document}